\section{Energy derivatives (forces, stress) and geometry optimization}

With self-consistent Kohn-Sham orbitals, densities and total energies
available, one of the primary tasks of electronic structure theory is
to obtain energy \emph{derivatives} with respect to the nuclear
coordinates $\boldR_\text{at}$. First derivatives (``forces'') allow
to find the optimum structure and molecular dynamics on the
Born-Oppenheimer surface. Based on the structure optimum, second
derivatives then enable the calculation of the (Born-Oppenheimer)
zero-point vibrational energy of the nuclei, and of vibrational
excitations (phonons).

The present section deals with the options available in FHI-aims for
the computation of forces in FHI-aims, and with algorithms related to
geometry optimization. Specifics regarding first-principles molecular
dynamics are given in Sec. \ref{Sec:MD}, and the computation of second
energy derivatives (vibrational frequencies, zero-point energies, and
oscillator strengths) by a finite-difference technique is covered in 
Sec. \ref{Sec:vib}. 

Before coming to the full set of related keywords, 
we give some basic comments below. Please also consider using (and 
adding to) the wiki, especially if you encounter any kind of behavior
that is obviously not intended. We would like to know about such
cases. 

\textbf{One ``most important'' rule:}

For efficient local structure optimizations, do not simply use
``tight'' settings from an arbitrary starting geometry and wait.

It is usually \emph{much} more effective to use ``light'' settings
first in a pre-relaxation step. Then, use ``tight'' settings only for
post-processing, e.g., based on the ``geometry.in.next\_step'' file 
that is now written out by the code by default.

\textbf{Basic handling:}

For most applications, simply specifying a convergence
criterion \keyword{sc\_accuracy\_forces} will be enough to switch on
the computation of forces itself. Reliable structure optimizations of
atomic coordinates can then be obtained using the
\keyword{relax\_geometry} with the \texttt{bfgs} algorithm. 

A note on naming conventions: \emph{The default version of BFGS calls
  the trust radius enhanced variant of BFGS, also available as
  \texttt{trm}. A standard BFGS variant (sub-keyword
  \texttt{bfgs\_textbook}) is also available as a fallback.} 

Beginning with code version ``012213'', we have also formalized the
handling of \emph{restarting} relaxations more clearly. If
\keyword{relax\_geometry} \texttt{bfgs} is used, a file
\texttt{geometry.in.next\_step} is written out by default which contains
both the geometry and the current estimate of the Hessian matrix of
the system that are used by the BFGS algorithm. For an organized
restart, simply copy this file to \texttt{geometry.in}, and the stored
geometry and Hessian will be used to reinitialize the BFGS algorithm
if the \texttt{control.in} file requests a structure relaxation.

There are also keywords \keyword{hessian\_to\_restart\_geometry} and
\keyword{write\_restart\_geometry} that can be set to ``.false.'' to
suppress the writing of the Hessian information, or even of the
\texttt{geometry.in.next\_step} file as a whole, if that information
is deliberately not wanted. 

\textbf{Unit cell shape optimization:}

In periodic calculations, you may also want to relax the unit cell
shape itself. In this case, an \emph{additional} keyword
\keyword{relax\_unit\_cell} must be set. Where it is available, an
analytically computed stress tensor will be used, thanks to work by
Viktor Atalla, Christian Carbogno, and Franz Knuth~\cite{Knuth2015}. For other
functionals, the unit cell itself can still be relaxed, but the stress
tensor must be computed numerically from finite differences of total
energies. 

The numerical finite-difference stress tensor is always available as a
fallback.

However, we also note that, in some cases, just optimizing the unit
cell shape blindly may not be what you want. For example, in
high-symmetry structures a fit to Murnaghan's Equation of state---see
the utility provided in the \emph{utilities} directory---will be more
accurate and give you more information about the system (bulk moduli
and, in case of more than one phase, also access to transition pressures).

\textbf{Hybrid functionals:}

For hybrid functionals, analytic energy gradients and the analytic stress 
tensor are only available together with the \keyword{RI\_method} 
\texttt{LVL\_fast} version of ``resolution of identity'' of the two-electron 
Coulomb operator. This is a localized, linear-scaling version of the more 
traditional ``RI-V'' method of quantum chemistry (see Ref. \cite{Ren12a} for a
comprehensive introduction of ``RI-V''). 

At present, geometry optimizations and molecular dynamics using 
hybrid functionals work, but \keyword{RI\_method} \texttt{LVL\_fast} must be 
used.

You may wish to check any resulting total energies by single point 
calculations (no relaxation) using \keyword{RI\_method} \texttt{V}
afterwards.

Note also that \keyword{RI\_method} \texttt{LVL\_fast} itself becomes 
more accurate as one goes to larger basis sets. One may, in fact, augment
only the auxiliary basis that is used to expand the Coulomb operator
in order to increase the accuracy of the Coulomb operator itself even when
a small orbital basis set is used. The keyword that allows one to do this
is \subkeyword{species}{for\_aux}. 

For perturbative methods (MP2 perturbation theory, RPA and beyond),
analytical gradients are not yet available. 

\newpage

\subsection*{Tags for \texttt{geometry.in}:}

\keydefinition{constrain\_relaxation}{geometry.in}
{
  \noindent 
  Usage: \keyword{constrain\_relaxation} \option{constraint} \\[1.0ex]
  Purpose: In \texttt{geometry.in}, fixes the position of the last
    specified \keyword{atom}/\keyword{lattice\_vector} in a structure optimization. \\[1.0ex]
  \option{constraint} is a string, indicating what exactly will be
    constrained. Default: \texttt{.false.} \\
}
Allows to relax only parts of a structure, while keeping the rest at
fixed positions. Currently, the following simple options for
  \texttt{constraint} are possible:  
\begin{itemize}
  \item \texttt{.true.}: All coordinates for this atom will be constrained. 
  \item \texttt{.false.}: The relaxation of this atom will not be constrained.
  \item \texttt{x}: The $x$ coordinate of this atom is not allowed to move.
  \item \texttt{y}: The $y$ coordinate of this atom is not allowed to move.
  \item \texttt{z}: The $z$ coordinate of this atom is not allowed to move.
\end{itemize}
\emph{Attention}: If you wish to constrain more than one coordinate, the required
constraints must be specified as separate lines, like this:
\begin{verbatim}
  atom 0. 0. 0. Fe
    constrain_relaxation x
    constrain_relaxation y
\end{verbatim}
In contrast, specifying two constraints in one line will \emph{not} work. The
second constraint would simply be ignored!

\keydefinition{hessian\_block}{geometry.in}
{
  \noindent 
  Usage: \keyword{hessian\_block} \option{i\_atom} \option{j\_atom} \option{block}\\[1.0ex]
  Purpose: In \texttt{geometry.in}, allows to specify a Hessian matrix
  explicitly, with one line for each 3$\times3$ block.  The option
  \option{block} consists of 9 numbers in column-first (Fortran) order.  The
  3$\times$3 block corresponding to \option{j\_atom}, \option{i\_atom} is
  initialized by the transposed of \option{block}. The Hessian matrix is input
  in units of eV/\AA$^2$.\\
} 

If at least one \keyword{hessian\_block} line is found in the file, the
Hessian is constructed using this mechanism.  So far there is no safe-guard
from overriding Hessian blocks with subsequent lines with equal
\option{i\_atom}, \option{j\_atom}.

There are two scripts in the \texttt{utilities} directory to automatically
construct such Hessian matrix approximations.  First,
\texttt{conversions/thess2aims.py} converts a Tinker generated Hessian matrix.
Second, \texttt{Lindh.py} constructs a general purpose model matrix
\cite{Lin95}.  Please note that the Lindh model matrix is now also directly
available with \keyword{init\_hess} \texttt{Lindh}.

\keydefinition{hessian\_block\_lv}{geometry.in}
{
  \noindent 
  Usage / purpose: Like \keyword{hessian\_block}, but for Hessian
  matrix elements between lattice vector degrees of freedom.\\
} 

\keydefinition{hessian\_block\_lv\_atom}{geometry.in}
{
  \noindent 
  Usage / purpose: Like \keyword{hessian\_block}, but for Hessian
  matrix elements between lattice vector degrees of freedom.\\
} 

\keydefinition{trust\_radius}{geometry.in}
{
  \noindent 
  Usage: \keyword{trust\_radius} \option{value} \\[1.0ex]
  Purpose: In \texttt{geometry.in}, allows to specify the initial
  trust radius value for the \texttt{trm} relaxation algorithm. \\ 
} 

This keyword is a significant exception -- it is the only
algorithmic keyword found in the \texttt{geometry.in}
file. Conceptually, it would belong into \texttt{control.in}, but
since it is written out as part of the \texttt{geometry.in.next\_step}
file which can be used to restart a structure relaxation, we keep it
here. 

The keyword specifies the initial value of the ``trust radius'' used
to limit structure relaxation steps determined by the \texttt{bfgs}
(synonymous with \texttt{trm}) relaxation algorithm.

The default is set to 0.2~{\AA}. It will be overridden by the default
value of \keyword{max\_atomic\_move}, which can be set in
\texttt{control.in}. 

\newpage

\subsection*{Tags for general section of \texttt{control.in}:}


\keydefinition{bfgs\_extrapolation\_cap}{control.in}
{
  \noindent
  Usage: \keyword{bfgs\_extrapolation\_cap} \option{value} \\[1.0ex]
  Purpose: If the \texttt{bfgs\_textbook} algorithm of \keyword{relax\_geometry}
    deviated from a line step of 1.0, sets a maximum for the line step
    during relaxation. \\[1.0ex] 
  \option{value} is a real number $\ge$1.0 (dimensionless), the
    maximum allowed step when a variable line step used with
    BFGS. Default: 4.0 \\
}
\option{value} is dimensionless, given in units specified by the BFGS
Hessian matrix. \option{value}=1.0 corresponds to the default
extrapolation of BFGS, assuming a perfectly quadratic energy surface.


\keydefinition{bfgs\_step}{control.in}
{  

  \noindent
  Usage: \keyword{bfgs\_step} \option{type} [\option{factor} \option{cap}] \\[1.0ex]
  Purpose: Can specify a BFGS extrapolation with an adaptive line
    step instead of the default linestep for purely quadratic
    extrapolation. \\[1.0ex]
  \option{type} is a string, either \texttt{quadratic} or
  \texttt{mix}. Default: \texttt{quadratic}. \\
  \option{factor} : If \option{type} is not \texttt{quadratic},
    specifies a mixing factor to interpolate between the line step
    actually used in the last relaxation step, and the \emph{a
    posteriori} estimated optimum line step. Default: 0.1 \\
  \option{cap} : If \option{type} is not \texttt{quadratic},
    specifies an upper limit for the allowed line step. \\
}

\keydefinition{calc\_analytical\_stress\_symmetrized}{control.in}
{
  \noindent
  Usage: \keyword{calc\_analytical\_stress\_symmetrized} \option{flag} \\[1.0ex]
  Purpose: If \texttt{.false.}, calculates all 9 components of the analytical
  stress tensor. If \texttt{.true.} calculates only the upper triangle (6
  components) of the tensor and copies the result to the lower triangle. \\[1.0ex] 
  \option{flag} is a logical string, either \texttt{.true.} or
    \texttt{.false.} Default: \texttt{.true.}\\
}
Generally, it is sufficient to calculate the upper triangle of the
tensor. This flag is mainly for debugging purposes.

\keydefinition{clean\_forces}{control.in}
{
  \noindent
  Usage: \keyword{clean\_forces} \option{type} \\[1.0ex]
  Purpose: Can remove small unitary force components (rotation and
    translation of the whole structure due to residual numerical
    noise) in relaxations. \\[1.0ex] 
  \option{type} is a string, specifying whether and how any overall
    rotations / translations are removed. \\
}
The default for \option{type} depends on the exact circumstances (see
below). The following choices exist: 
\begin{itemize}
  \item \texttt{none} : No removal of residual rotations and
    translations. This is the default if any external embedding fields
    or charges are specified in \texttt{geometry.in}.
  \item \texttt{sayvetz} : Non-periodic structures: Removal of
    rotations and translations by a formal projection
    \cite{Eckart35,Sayvetz39}. In \emph{periodic} systems,
    only translations are removed.
  \item \texttt{fixed\_plane} : \emph{experimental} Simple alternative
    algorithm by constraining three atoms into a plane (implicitly
    constraining all others). 
\end{itemize}

\keydefinition{compute\_analytical\_stress}{control.in}
{
  \noindent
  Usage: \keyword{compute\_analytical\_stress} \option{flag} \\[1.0ex]
  Purpose: If \texttt{.true.}, switches on the computation of
  the analytical stress tensor. \\[1.0ex] 
  \option{flag} is a logical string, either \texttt{.true.} or
    \texttt{.false.} \\
}
Default: \texttt{.true.} if a unit cell relaxation was requested and computation
is possible. Otherwise, \texttt{.false.}

This flag allows to request an explicit analytical stress tensor computation
for an otherwise explicit single-point calculation.

The calculation of the analytical stress is limited to LDA, GGA, Meta-GGA and 
hybrid functionals and is not possible with \keyword{load\_balancing}. The vdW
correction based on Hirshfeld partitioning (\keyword{vdw\_correction\_hirshfeld})
is included into the analytical stress tensor.

\keydefinition{compute\_forces}{control.in}
{
  \noindent
  Usage: \keyword{compute\_forces} \option{flag} \\[1.0ex]
  Purpose: If \texttt{.true.}, switches on the computation of
    forces. \\[1.0ex] 
  \option{flag} is a logical string, either \texttt{.true.} or
    \texttt{.false.} \\
} 
Default: \texttt{.true.} if a geometry optimization or
molecular dynamics run was requested, or if the 
\keyword{sc\_accuracy\_forces} convergence criterion was
set. Otherwise, \texttt{.false.} 

This flag allows to request an explicit force computation for an
otherwise explicit single-point calculation. This is necessary for use
with external tools that require forces, such as a finite-difference
calculation of vibrational frequencies (see
Sec. \ref{Sec:vib}) or a transition state search (see
Sec. \ref{appendix_NEB}). In these cases, keyword
\keyword{final\_forces\_cleaned} should also be set.

\keydefinition{compute\_numerical\_stress}{control.in}
{
  \noindent
  Usage: \keyword{compute\_numerical\_stress} \option{flag} \\[1.0ex]
  Purpose: If \texttt{.true.}, switches on the computation of
  the numerical stress tensor based on central finite differences. \\[1.0ex] 
  \option{flag} is a logical string, either \texttt{.true.} or
    \texttt{.false.} \\
}
Default: \texttt{.true.} if a unit cell relaxation was requested and the
computation of the analytical stress is not possible. Otherwise,
\texttt{.false.}

If not further specified (by  \keyword{delta\_numerical\_stress}) the default
value for the scaling factor delta is set to $10^{-4}$.

\keydefinition{delta\_numerical\_stress}{control.in}
{
  \noindent
  Usage: \keyword{delta\_numerical\_stress} \option{value}                        \\[1.0ex]
  Purpose: Specifies the scaling factor delta in the computation
           of the numerical stress tensor (\keyword{compute\_numerical\_stress}). \\[1.0ex] 
  \option{value} is a dimensionless real number $> 0$. Default: $10^{-4}$.
}

\keydefinition{energy\_tolerance}{control.in}
{
  \noindent
  Usage: \keyword{energy\_tolerance} \option{tolerance} \\[1.0ex]
  Purpose: Sets the energy amount by which a relaxation step can move
    upwards and is still accepted. \\[1.0ex]
  \option{tolerance} is a small positive real number, in eV. Default:
    2$\cdot$10$^{-4}$~eV. \\
}
Small upward steps during relaxation may occur as a result of a
slightly mis-guessed \texttt{bfgs\_textbook} Hessian matrix somewhere
along the path, or
as a result of some residual numerical noise that leads to a
discrepancy between energies and forces. In the present code version,
such noise is always safely below the default \option{energy\_tolerance} for
reasonable settings. If exceeded, this may point to a different
problem, simply increasing \option{tolerance} is not, in general, a
good idea.  For the \texttt{trm} optimizer, also see
\keyword{harmonic\_length\_scale}.

\keydefinition{external\_force}{geometry.in}
{
  \noindent
  Usage: \keyword{external\_force} \option{x} \option{y} \option{z} \\[1.0ex]
  Purpose: \emph{Experimental} -- Applies an external force to the atom previous to this keyword.\\[1.0ex]
  \option{x},\option{y},\option{z} are the force components in eV/\AA\ applied to the atom in $x$, $y$, $z$ direction.\\
}
When an external force is applied it is necessary to contstrain the relaxation of at least on other atom to avoid a constant shift of the geometry. Also the value has to be reasonably chosen. Tear apart geometries can result in very flat energy landscapes, which take a large amount of time to optimize. Typically this feature should be used in cases where a small external force, e.g. a STM-tip is applied on an atomic layer and the geomtry response of this external force is of interest.

\emph{This feature is experimental since no extensive testing was done for it.}


\keydefinition{external\_pressure}{control.in}
{
  \noindent
  Usage: \keyword{external\_pressure} \option{value} \\[1.0ex]
  Purpose: \emph{Experimental} -- Applies external pressure to the unit cell.\\[1.0ex]
  \option{value} is the pressure in eV/\AA$^3$ applied to the unit cell.\\
}
In the periodic case, it is possible to apply hydrostatic pressure to the unit cell.
To actually see the effect of the external pressure, a unit cell relaxation is required
(see \keyword{relax\_geometry} and \keyword{relax\_unit\_cell}). The crystal is then
relaxed with the external pressure added to the stress tensor.

\emph{This feature is experimental since no extensive testing was done for it.}

\keydefinition{final\_forces\_cleaned}{control.in}
{
  \noindent
  Usage: \keyword{final\_forces\_cleaned} \option{flag} \\[1.0ex]
  Purpose: Decides whether spurious unitary transformations of the
    complete system (translations and rotations) are removed before
    the final output. \\[1.0ex]
  \option{flag} is a logical string, either \texttt{.true.} or
    \texttt{.false.} Default: \texttt{.true.} \\
}
This option affects directly the long-format (15 decimal) output of
total energies and forces at the end of the s.c.f. cycle in the standard
output file. If \option{flag} is \texttt{.true.}, the final output
forces are ``cleaned'' using the \texttt{sayvetz}
  \cite{Eckart35,Sayvetz39} mechanism of keyword  
\keyword{clean\_forces} (removal of translations and rotations for
cluster geometries; only translations removed for periodic systems). 

\keyword{final\_forces\_cleaned} \texttt{.true.} should be set for use
with external tools that require forces, such as a finite-difference
calculation of vibrational frequencies (see
Sec. \ref{Sec:vib}) or a transition state search (see 
Sec. \ref{appendix_NEB}). 

\keydefinition{harmonic\_length\_scale}{control.in}
{
  \noindent
  Usage: \keyword{harmonic\_length\_scale} \option{length} \\
  Purpose: The \texttt{trm}/\texttt{bfgs} algorithm of
  \keyword{relax\_geometry} judges a step by its energy gain.  Usually, one
  simply uses the energy difference.  For very short steps and rather light
  grids, however, it turns out that the qualitiy of the energy is inferior to
  the quality of the forces.  For steps shorter than \option{length}, do not
  look at the energy but use the harmonic approximation $- \Delta
  \tilde E = (\bm X_2-\bm X_1)\cdot(\bm F_2+\bm F_1)/2$ as an estimate for the
  energy gain.  If this procedure willfully accepts
  a step which increases the energy by more than \keyword{energy\_tolerance},
  the code stops to warn the user about the inconsistency between energy
  functional and forces.
  \\
  \option{length} is a length scale in \AA. Default: \texttt{0.025} \\
}
Effectively, this flag switches from a real energy minimizer to a search for a
stable zero of the force field for short step lengths.


\keydefinition{hessian\_to\_restart\_geometry}{control.in}
{
  \noindent
  Usage: \keyword{hessian\_to\_restart\_geometry} \option{flag} \\
  Purpose: Exports the current approximation to
  the Hessian matrix to \option{geometry.in.next\_step} during a 
  relaxation restart using \keyword{hessian\_block}. \\ 
  \option{flag} is a logical string, either \texttt{.true.} or
    \texttt{.false.} Default: \texttt{.true.} \\
}

Note: The \texttt{geometry.in.next\_step} file is written out by
default when the \keyword{relax\_geometry} \texttt{bfgs} algorithm
is used. For the older \texttt{bfgs\_textbook} algorithm, this file is
only written when combined with the \keyword{restart\_relaxations}
keyword. 

\keydefinition{init\_hess\_lv\_diag}{control.in}
{
  \noindent
  Usage: \keyword{init\_hess\_lv\_diag} \option{value} \\[1.0ex]
  Purpose: In a geometry relaxation with a unit cell optimization,
    allows to specify the initial Hessian matrix elements used to
    estimate relaxation steps that are associated with the lattice
    vector degrees of freedom.
  \\[1.0ex]
  \option{length} is a length scale in eV/{\AA$^2$}. Default: 25
  ev/{\AA$^2$} \\ 
}
See the \keyword{init\_hess} keyword for more information on the
initial Hessian matrix used during a structure optimization.

\keydefinition{init\_hess}{control.in}
{
  \noindent
  Usage: \keyword{init\_hess} \option{diag} [\option{value}]
  \emph{or:} \keyword{init\_hess} \option{Lindh} [\option{add-value}
  [\option{thres}]]
  \\[1.0ex]
  Purpose: Defines the initial Hessian matrix that is used by the
  structure relaxation algorithms \texttt{bfgs\_textbook} and
  \texttt{bfgs} (the latter is synonymous with \texttt{trm}) of the
  \keyword{relax\_geometry} keyword.
  \\[1.0ex]
  Default: \texttt{init\_hess Lindh 2.} unless a specific Hessian is
  given in the \texttt{geometry.in} file. \\ 
}

If no explicit initial Hessian matrix is given in the geometry.in file
or specified using the older \keyword{restart\_relaxations}
infrastructure, the keyword \keyword{init\_hess} can be used to
specify the initial Hessian matrix used in a structure optimization. 

With the option \option{diag}, a diagonal initial Hessian matrix is
assumed. Then, the number given by the \option{value} option sets the
diagonal elements between all atomic coordinates directly. The default
is 25\,eV/\AA$^2$. Larger values lead to a more conservative
start, smaller values lead to more aggressive initial relaxation steps.

With \option{Lindh} the Lindh model matrix \cite{Lin95} is used to
initialize the Hessian between all \emph{atomic} coordinates (usually
a very efficient guess). For  stability reasons, \option{add-value}
(in eV/\AA$^2$, defaults to 2.0\,eV/\AA$^2$) is added to all matrix 
elements on the diagonal. The parameter \option{thres} (default: 
15.0) can be used to specify the accuracy of the Lindh matrix; only
terms estimated to be larger than $e^{-\text{\option{thres}}}$ are
taken into account.   

If a unit cell relaxation is additionally requested, all the Hessian
rows and columns associated with the lattice vector degrees of freedom
are set to be diagonal, both for the \option{diag} and \option{Lindh}
versions. The default diagonal element is 25\,eV/\AA$^2$. Its value
can be modified using the \keyword{init\_hess\_lv\_diag} keyword. If
you are planning a large number of complex unit cell optimizations of
a similar type, we do recommend to check whether this default value is
any good. 

If the initial Hessian is specified explicitly by
\keyword{hessian\_block} in \texttt{geometry.in}, this explicit
hessian overrides any information requested by the
\keyword{init\_hess} keyword.


\keydefinition{line\_search\_tol}{control.in}
{
  \noindent
  Usage: \keyword{line\_search\_tol} \option{tolerance} \\[1.0ex]
  Purpose: During a \keyword{relax\_geometry} \texttt{bfgs\_textbook} geometry
    relaxation, rejects a
    relaxation step if the \emph{a posteriori} estimate for the optimal
    relaxation step deviates by more than a given tolerance from the
    actually performed line step. \\[1.0ex]
  \option{tolerance} is a real positive relative tolerance
    (dimensionless in units of the BFGS Hessian). Default:
    3.0 \\ 
}
If keyword \keyword{relax\_geometry} \texttt{bfgs\_textbook} is used,
\keyword{line\_search\_tol} will reject a relaxation step based on an
\emph{a posteriori} estimate of the true optimum step length (see 
Ref. \cite{Blum08} for details). In short, a step of length $l$ in
dimensionless units (units relative to the BFGS Hessian matrix) is
rejected if 
\begin{equation}
  |l_\text{\emph{a posteriori}} - l_\text{actual}|/l_\text{actual} \ge
   \mathtt{tolerance} \, .
\end{equation}

\keydefinition{line\_step\_reduce}{control.in}
{
  \noindent
  Usage: \keyword{line\_step\_reduce} \option{value} \\[1.0ex]
  Purpose: Determines the behaviour of FHI-aims if a
    \keyword{relax\_geometry} \texttt{bfgs\_textbook} relaxation step was
    rejected. \\[1.0ex]
  \option{value} is either a string descriptor, or a numerical
    value. Default: \texttt{automatic} . \\
}
During a structure optimization using the \keyword{relax\_geometry}
\texttt{bfgs\_textbook} algorithm, FHI-aims checks if a performed relaxation
step actually reduced the total energy compared to the
previous step. The allowed upward tolerance is set by keyword
\keyword{energy\_tolerance}. If an upward step is encountered,
FHI-aims first attempts to repeat the step with a smaller step
length. If this smaller step also leads upwards, FHI-aims assumes
that the BFGS Hessian has exceeded its validity limits, and
reinitializes the Hessian before continuing. 

If \texttt{value} is set to \texttt{automatic}, the BFGS step is
repeated with an \emph{a posteriori} estimated length $l$ given by the
previous (failed) step (see Ref. \cite{Blum08}).

Alternatively, \texttt{value} can be set to a real positive number $<$1,
corresponding to the default factor by which the previous step length
is reduced. 

\keydefinition{max\_atomic\_move}{control.in}
{
  \noindent
  Usage: \keyword{max\_atomic\_move} \option{value} \\[1.0ex]
  Purpose: Maximum allowed step length taken during relaxation. \\[1.0ex]
  \option{value} is a real positive upper bound for the maximum allowed
    change in single atomic coordinate, in {\AA}. Default: 0.2 {\AA}. \\
}
If the \texttt{bfgs\_textbook}-predicted change in an atomic coordinate exceeds
\keyword{max\_atomic\_move}, the length of the entire step (all
coordinates) will be scaled down to not exceed the maximum allowed
displacement.

\keydefinition{max\_relaxation\_steps}{control.in}
{
  \noindent
  Usage: \keyword{max\_relaxation\_steps} \option{number} \\[1.0ex]
  Purpose: A structure optimization will be aborted after exceeding a
    prescribed maximum number of steps. \\[1.0ex]
    \option{number} is the prescribed maximum step number. Default:
    1000 . \\
}

\keydefinition{min\_line\_step}{control.in}
{
  \noindent
  Usage: \keyword{min\_line\_step} \option{value} \\[1.0ex]
  Purpose: Sets the minimum possible (dimensionless) BFGS linestep
    before the BFGS Hessian is reinitialized \\[1.0ex]
  \option{value}: A real positive number (between 0 and 1). Default:
    0.1 \\
}
See keyword \keyword{line\_step\_reduce}. If, after a rejected
\texttt{bfgs\_textbook} relaxation step, the reduced, dimensionless BFGS line
step decreases below \keyword{min\_line\_step}, the BFGS Hessian
matrix estimate is considered inappropriate for the present point in
space, and the Hessian is reinitialized.

\keydefinition{numerical\_stress\_save\_scf}{control.in}
{
  \noindent
  Usage: \keyword{numerical\_stress\_save\_scf} \option{flag} \\[1.0ex]
  Purpose: Controls if \keyword{constrain\_relaxation}
  directives are used to determine implicitly if a component of the numerical
  stress has to be calculated. This greatly accelerates unit cells
  with high symmetries (e.g. orthorhombic).  \\[1.0ex]
 \option{flag} is a logical string, either \texttt{.true.} or
    \texttt{.false.} Default: \texttt{.true.} \\
}

\keydefinition{orthonormalize\_eigenvectors}{control.in}
{
  \noindent
  Usage: \keyword{orthonormalize\_eigenvectors} \option{flag} \\[1.0ex]
  Purpose: Specifies whether or not the wave function coefficients
    from the previous geometry will be re-orthonormalized before
    initializing a new relaxation step. \\[1.0ex]
  \option{flag} is a logical string, either \texttt{.true.} or
    \texttt{.false.} Default: \texttt{.true.} \\
}
The \keyword{orthonormalize\_eigenvectors} keyword allows to
reorthnormalize the converged self-consistent Kohn-Sham orbitals
$c_{jl}$ after a relaxation step. These are then used to
reinitialize the electron density for the next relaxation step.

Due to the change in atomic positions, the wave function coefficients
$c_{jl}$ for the earlier geometry are no longer orthonormal after the
relaxation step. The consequence is an initial electron density which
no longer satisfies the correct electron count (i.e., the system may
appear to be charged immediately after a relaxation step, although a
neutral system was requested). In principle, this does not matter for
the outcome of a calculation, since the self-consistent solution will
be independent of the starting point. In many cases, there is no clear
benefit in terms of the s.c.f. convergence duration from
orthonormalizing the $c_{jl}$ prior to the reinitialization; however, some
cases with unstable s.c.f. convergence may benefit significantly.

\keydefinition{relax\_geometry}{control.in}
{
  \noindent
  Usage: relax\_geometry \option{type} \option{tolerance} \\[1.0ex]
  %
  Purpose: Specifies if a structure optimization (geometry relaxation)
    is requested, and which. \\[1.0ex]
  %
  \option{type} specifies the type of requested structure
    optimization. Default: \texttt{none}. \\
  \option{tolerance}: Specifies the maximum residual
    force component per atom (in eV/{\AA}) below which the geometry
    relaxation is considered converged. \\
}
Finds the nearest minimum of the Born-Oppenheimer potential energy
surface for the nuclei.

For periodic calculations: If you are looking to relax not just atomic
coordinates but also the unit cell shape (lattice vectors), you do
need to specify an additional keyword: \keyword{relax\_unit\_cell}.

The presently supported options for \option{type} are \texttt{none},
\texttt{bfgs\_textbook} and \texttt{trm} (which is now synonymous with
\texttt{bfgs}).
\begin{itemize}
  \item \texttt{bfgs} or \texttt{trm} uses a trust radius method enhanced
    version of the Broyden-Fletcher-Goldfarb-Shanno (BFGS) optimization
    algorithm (see Ref. \cite{Nocedal06-numopt}, which was the basis for
    J\"urgen Wieferink's code effort in this area). In our tests, this version
    appears to give the fastest convergence reliably.
  \item \texttt{bfgs\_textbook} invokes a Cartesian BFGS optimization algorithm
    that adheres strictly to the textbook \cite{NumericalRecipes}; we note that
    our coded version is based on the concise description found at Wikipedia.
    See Ref. \cite{Blum08} for some additional details of our version.
\end{itemize}

The relaxation algorithm can be greatly sped up by using a somewhat
intelligent guess for the Hessian matrix used in the initial step. 
By default, FHI-aims now sets the general purpose model matrix due to Lindh
and coworkers \cite{Lin95} with a slight modification. For more information
see the \keyword{init\_hess} keyword above.

A reliable force convergence criterion \option{tolerance} for most
structures is 10$^{-2}$ eV/{\AA} or 5$\cdot$10$^{-3}$ eV/{\AA}.  

Going to a much smaller \texttt{tolerance} value can be necessary for
some specific purposes, but do it with great care. If
\texttt{tolerance} is set to a too small value by default, 80\% or
more of your CPU time may be spent groping around in the last meV of
the structure optimization.

Use the 'get\_relaxation\_info.pl' utility script to monitor the
progress of an ongoing or finished relaxation run.

To stop an ongoing structure relaxation in an organized way, create
the \keyword{abort\_opt} file in the respective directory.

For vibrations of large molecules or phonons in solids from finite
differences, tighter settings of the \texttt{tolerance} parameter may
be required, even down to 10$^{-4}$ eV/{\AA}. In that case, however,
do not forget that rather tight s.c.f. convergence accuracy settings
may also be required to get accurate gradients in the first place. 

In other words, use the \texttt{tolerance} criterion for a structure
relaxation run wisely -- decide what is the physical quantity you are
actually interested in, and then check which value of the
\texttt{tolerance} criterion is safe but still efficient.


\keydefinition{relax\_unit\_cell}{control.in}
{
  \noindent
  Usage: \keyword{relax\_unit\_cell} \option{type}  \\[1.0ex] 
  Purpose: Relaxes unit cell (lattice vectors) using the structure 
           optimization as specified in \keyword{relax\_geometry}. \\[1.0ex]
  \option{type} specifies the type of requested unit cell 
    optimization. Presently supported options: \texttt{none},
    \texttt{full}, \texttt{fixed\_angles}. Default: \texttt{none}\\
}
Allows to optimize the lattice vectors of a periodic calculation, in
addition to the normal atomic coordinates inside the unit cell. This
keyword is not on by default, as automatically optimizing the unit
cell of (say) a surface calculation could do a lot of unintended
harm. Possible settings:
\begin{itemize}
  \item \texttt{none} : Unit cell will be kept fixed, no
    optimization. 
  \item \texttt{full} : All \keyword{lattice\_vector} degrees of
    freedom will be relaxed, except those affectd by explicit
    constraints.
  \item \texttt{fixed\_angles} : All angles between lattice vectors
    will be constrained (kept fixed), only the lengths of each lattice
    vector are varied. (This option used to called \texttt{shape}, but
    that is a misunderstandable term. The \texttt{shape} term will be
    removed in the future to avoid confusion.)
\end{itemize}
This keyword should be used only together with \keyword{relax\_geometry}.
Individual lattice vectors or its components can be constrained by using 
\keyword{constrain\_relaxation}.

If the computation of the analytical stress is possible regarding the chosen
computational settings, the analytical stress is used for the unit cell
relaxation. Otherwise, the numerical stress is used. With
\keyword{stress\_for\_relaxation}, one can explicitly choose either numerical or
analytical stress for the unit cell relaxation.

If a unit cell relaxation produces strange results with the analytical
stress, here are some potential remedies:
\begin{enumerate}
  \item A very possible reason may be because the integration grids
    are not dense enough. This could especially well be the case for
    ``light'' settings. One remedy is to just use the integration
    grids from ``tight'' and the basis functions from ``light''. 
  \item Another possible remedy is to switch the way the integration
    weights are calculated to a slightly slower, non-default
    version. E.g., change the \keyword{partition\_type} to a spherical
    one like \texttt{rho\_r2}. 
  \item Finally, you may wish to set the convergence of the analytical
    stress with \keyword{sc\_accuracy\_stress} to an explicit, final
    value. \textbf{Only ever set this keyword for test purposes,
    though, not routinely. The cost for too many analytical stress
    evaluations can be disproportionately large.}
\end{enumerate}

\keydefinition{stress\_for\_relaxation}{control.in}
{
  \noindent
  Usage: \keyword{stress\_for\_relaxation} \option{type} \\[1.0ex]
  Purpose: Use either numerical or analytical stress for unit cell
  relaxations.\\[1.0ex] 
  \option{type} can be either \texttt{numerical} or
    \texttt{analytical}. Default: Chosen automatically based on computational
    settings. \\
}
To perform an actual unit cell relaxation, one has to set
\keyword{relax\_unit\_cell}. If one chooses \texttt{analytical} but the 
computation of the analytical stress is not possible, FHI-aims will abort.

\keydefinition{restart\_relaxations}{control.in}
{
  \noindent
  Usage: \keyword{restart\_relaxations} \option{flag}
    \\[1.0ex]
  Purpose: To save all necessary data to restart a \texttt{bfgs\_textbook}
    relaxation without needing to recompute the relaxation path in
    DFT. \\[1.0ex] 
 \option{flag} is a logical string, either \texttt{.true.} of
 \texttt{.false.} Default: \texttt{.false.} \\
}

Note: The \keyword{restart\_relaxations} keyword will work with both
the \texttt{bfgs} (\texttt{trm}) or the \texttt{bfgs\_textbook}
versions of the \keyword{relax\_geometry} keyword. However, it is only
needed for the \texttt{bfgs\_textbook} version, if a restart is
desired. For the normal (and faster) \texttt{bfgs} (\texttt{trm})
version, all that is needed is the information contained in
\texttt{geometry.in.next\_step}. 

In order to restart
from an already known estimated Hessian, rather from unity (the
default in the absence of any previous knowledge), the Hessian can be
rebuilt from a series of stored atomic cordinates and forces.  

This information is exported to a separate file
\texttt{relaxation\_restart\_file.FHIaims}, which will be read and
evaluated by a newly beginning FHI-aims calculation if the
\keyword{restart\_relaxations} keyword is set to true.

The flag also needs a \texttt{geometry.in} file that is an exact copy
of the \texttt{geometry.in.next\_step} file, which is updated during
the structure optimization.

See also \keyword{restart} for keeping the wave functions as well. 

\keydefinition{write\_restart\_geometry}{control.in}
{
  \noindent
  Usage: \keyword{write\_restart\_geometry} \option{flag} \\
  Purpose: During a structure optimization, exports the current
  geometrys and approximation to the Hessian matrix to a file 
  \option{geometry.in.next\_step}. \\ 
  \option{flag} is a logical string, either \texttt{.true.} or
    \texttt{.false.} Default: \texttt{.true.} \\
}

Note: The \texttt{geometry.in.next\_step} file is written out by
default when the \keyword{relax\_geometry} \texttt{bfgs} algorithm
is used. For the older \texttt{bfgs\_textbook} algorithm, this file is
only written when combined with the \keyword{restart\_relaxations}
keyword. 

